%\documentclass[handout]{beamer}
\documentclass{beamer}
 
\usetheme[numbering = fraction, progressbar = none, background = light, sectionpage = progressbar]{metropolis}
\usepackage{amsmath}
\usepackage{synttree}
\usepackage{tabu}
\usepackage{graphicx}

\title{Lecture 3 Review Quiz}
\author{Mallick Hossain}
\date{}
\institute{University of Pennsylvania}

\begin{document}

%%%%%%%%%%%%%%%%%%%%%%%%%%%%%%%%%%%%%%%%
\begin{frame}
	\titlepage 
\end{frame} 

%%%%%%%%%%%%%%%%%%%%%%%%%%%%%%%%%%%%%%%%
\begin{frame}
\frametitle{Class Logistics}
    \begin{enumerate}[<+- | alert@+>]
        \item Have you started studying for the midterm?
            \begin{itemize}
                \item The correct answer is ``Yes''
            \end{itemize}
       \item When is it?
       	\begin{itemize}
       		\item Next Wednesday
       	\end{itemize}
        \item Have you started thinking about your R project?
            \begin{itemize}
                \item The correct answer is ``Yes'' \only<7>{\alert{but I will also accept ``Kinda, I looked at some of the links you suggested''}}
            \end{itemize}
    \end{enumerate}
\end{frame} 

%%%%%%%%%%%%%%%%%%%%%%%%%%%%%%%%%%%%%%%%
\begin{frame}
\frametitle{Test Questions}
    \begin{enumerate}[<+- | alert@+>]
        \item Will we get a formula sheet?
            \begin{itemize}
                \item No. If there's a formula that does not make sense, ask and we can remedy that.
            \end{itemize}
       \item Will R be tested?
       	\begin{itemize}
       		\item Yes, but I will not ask you to write code on the exam. 
       	\end{itemize}
        \item What will be tested?
            \begin{itemize}
                \item Everything through the end of today's lecture
            \end{itemize}
        \item Will we still have class before/after the exam?
        		\begin{itemize}
        			\item Yes. The material we cover will depend upon how far we get today. 
        		\end{itemize}
        	\item Will you post practice exams?
        		\begin{itemize}
        			\item Yes
        		\end{itemize}
    \end{enumerate}
\end{frame} 

%%%%%%%%%%%%%%%%%%%%%%%%%%%%%%%%%%%%%%%%
\begin{frame}
\frametitle{Homework-Inspired Test Questions}
    \begin{enumerate}[<+- | alert@+>]
        \item Do my answers have to be organized?
        	\begin{itemize}
        		\item Yes. The easier it is for me to understand your reasoning and thought process, the easier it will be to give you points.
        	\end{itemize}
        	\item Does notation matter?
        		\begin{itemize}
        			\item Yes. $\mu$ and $\sigma$ are parameters while $\bar{x}$ and $s_x$ are statistics. This will become more important as we progress through the course.
        		\end{itemize}
        	\item Can I give examples instead of proofs?
        		\begin{itemize}
        			\item No. Examples can help frame your intuition, but the proof is what establishes the facts.
        		\end{itemize}
        	\item What other things do I need to be careful of?
        		\begin{itemize}
        			\item Equal signs. If you use equal signs, make sure the two sides are equal.
        			\item Proving things are equal. Start with what you know, not what you're trying to prove.
        		\end{itemize}
    \end{enumerate}
\end{frame} 

%%%%%%%%%%%%%%%%%%%%%%%%%%%%%%%%%%%%%%%%
\begin{frame}
\frametitle{Real Questions}
    \begin{enumerate}[<+- | alert@+>]
		\item What is the difference between a ratio and an interval variable?
			\begin{itemize}
				\item An interval variable only has meaningful differences, but no natural zero while a ratio variable has meaningful differences and ratios as well as a natural zero.
			\end{itemize}			        
        \item What is the gambler's fallacy?
            \begin{itemize}
                \item The belief that short-run probabilities should reflect long-run probabilities.
            \end{itemize}
        \item What is the formula for $P(A \cap B)$?
        	\begin{itemize}
        		\item $P(A) + P(B) - P(A \cup B)$
        	\end{itemize}
        \item How many different Chipotle burritos can I make? There are 5 meats, 2 rice, 2 beans, and 4 salsas. I can only pick one of each ingredient.
        	\begin{itemize}
        		\item 5 * 2 * 2 * 4 = 80 burritos! 
        	\end{itemize}
     \item Ready for the next lecture?	
    \end{enumerate}
\end{frame} 

%%%%%%%%%%%%%%%%%%%%%%%%%%%%%%%%%%%%%%%%
\begin{frame}
\frametitle{Real Questions}
    \centering
	   \includegraphics[scale = 0.5]{letsDoThis3.jpg}
\end{frame} 

\end{document}